\section{Multisensor Payload Architecture}

MANAR incorporates a software model representing eleven controllable payload categories for victim search, identification, and signaling.

\subsection{Payload Categories}
The prototype models eleven controllable component categories in \texttt{components.cpp}. The descriptions below indicate their intended subsystem roles; sensing and processing functionality is not yet implemented:
\begin{enumerate}
    \item \textbf{Thermal LWIR Camera:} Body heat detection.
    \item \textbf{Day RGB Camera:} Visual inspection and terrain validation.
    \item \textbf{Low-Light / IR Camera:} Low ambient light sensitivity.
    \item \textbf{FMCW Radar:} Range, velocity, and micro-motion detection.
    \item \textbf{Passive RF Sensor:} Mobile signal uplink detection.
    \item \textbf{Public Address Emitter:} Audible instruction broadcast.
    \item \textbf{Microphone Array:} Acoustic monitoring for vocal signals.
    \item \textbf{Amber Beacon:} Visual alerting for ground teams.
    \item \textbf{White Guidance Strobe:} Visual orientation strobe.
    \item \textbf{Downward Spotlight:} Ground illumination.
    \item \textbf{Smoke Marker Canister:} Visual smoke signaling.
\end{enumerate}

\subsection{Candidate Hardware Mapping}
Table~\ref{tab:hardware_mapping} maps software payload categories to candidate hardware under evaluation. Rows marked TBD indicate categories where hardware selection has not been evaluated.

\begin{table}[t]
\caption{Candidate Payload Hardware (Non-Final)}
\label{tab:hardware_mapping}
\centering
\footnotesize
\begin{tabular}{@{}ll@{}}
\toprule
\textbf{Category} & \textbf{Candidate Hardware} \\
\midrule
Thermal LWIR & FLIR Boson+ 640 \\
Day RGB & Sony FCB-EV9500L \\
Low-Light IR & SiOnyx Aurora \\
FMCW Radar & TBD (band not finalized) \\
Passive RF & TBD \\
PA Emitter & TBD \\
Microphone Array & ReSpeaker v2.0 \\
Visual Alerting & TBD \\
Smoke Marker & TBD \\
\bottomrule
\end{tabular}
\end{table}

\subsection{Perception Pipeline Integration}
In the current prototype, payload component states can be commanded individually via \texttt{CHANGE\_COMPONENTS} or toggled automatically during mission hold states. Automated sensor fusion algorithms and ML perception models are reserved for future work.
